\documentclass[10pt,a4paper]{amsart}
\usepackage{amsmath,amssymb,amsthm,graphicx}
\usepackage{verbatim}
\newenvironment{reference}{\comment}{\endcomment}
%\newenvironment{reference}{}{}
\newenvironment{slogan}{\comment}{\endcomment}
\newenvironment{history}{\comment}{\endcomment}

% For commutative diagrams we use Xy-pic
\usepackage[all]{xy}

% We use 2cell for 2-commutative diagrams.
\xyoption{2cell}
% Theorem environments.
%
\theoremstyle{plain}
\newtheorem{theorem}[subsection]{Theorem}
\newtheorem{proposition}[subsection]{Proposition}
\newtheorem{lemma}[subsection]{Lemma}

\theoremstyle{definition}
\newtheorem{definition}[subsection]{Definition}
\newtheorem{example}[subsection]{Example}
\newtheorem{exercise}[subsection]{Exercise}
\newtheorem{situation}[subsection]{Situation}

\theoremstyle{remark}
\newtheorem{remark}[subsection]{Remark}
\newtheorem{remarks}[subsection]{Remarks}

\numberwithin{equation}{subsection}

% Macros
%
\def\lim{\mathop{\mathrm{lim}}\nolimits}
\def\colim{\mathop{\mathrm{colim}}\nolimits}
\def\Spec{\mathop{\mathrm{Spec}}}
\def\Hom{\mathop{\mathrm{Hom}}\nolimits}
\def\Ext{\mathop{\mathrm{Ext}}\nolimits}
\def\SheafHom{\mathop{\mathcal{H}\!\mathit{om}}\nolimits}
\def\SheafExt{\mathop{\mathcal{E}\!\mathit{xt}}\nolimits}
\def\Sch{\mathit{Sch}}
\def\Mor{\mathop{\mathrm{Mor}}\nolimits}
\def\Ob{\mathop{\mathrm{Ob}}\nolimits}
\def\Sh{\mathop{\mathit{Sh}}\nolimits}
\def\NL{\mathop{N\!L}\nolimits}
\def\CH{\mathop{\mathrm{CH}}\nolimits}
\def\proetale{{pro\text{-}\acute{e}tale}}
\def\etale{{\acute{e}tale}}
\def\QCoh{\mathit{QCoh}}
\def\Ker{\mathop{\mathrm{Ker}}}
\def\Im{\mathop{\mathrm{Im}}}
\def\Coker{\mathop{\mathrm{Coker}}}
\def\Coim{\mathop{\mathrm{Coim}}}

% Boxtimes
%
\DeclareMathSymbol{\boxtimes}{\mathbin}{AMSa}{"02}

%
% Macros for moduli stacks/spaces
%
\def\QCohstack{\mathcal{QC}\!\mathit{oh}}
\def\Cohstack{\mathcal{C}\!\mathit{oh}}
\def\Spacesstack{\mathcal{S}\!\mathit{paces}}
\def\Quotfunctor{\mathrm{Quot}}
\def\Hilbfunctor{\mathrm{Hilb}}
\def\Curvesstack{\mathcal{C}\!\mathit{urves}}
\def\Polarizedstack{\mathcal{P}\!\mathit{olarized}}
\def\Complexesstack{\mathcal{C}\!\mathit{omplexes}}
% \Pic is the operator that assigns to X its picard group, usage \Pic(X)
% \Picardstack_{X/B} denotes the Picard stack of X over B
% \Picardfunctor_{X/B} denotes the Picard functor of X over B
\def\Pic{\mathop{\mathrm{Pic}}\nolimits}
\def\Picardstack{\mathcal{P}\!\mathit{ic}}
\def\Picardfunctor{\mathrm{Pic}}
\def\Deformationcategory{\mathcal{D}\!\mathit{ef}}
\usepackage{cleveref}
\crefname{originallemma}{original lemma}{original lemmas}
\Crefname{originallemma}{Original lemma}{Original lemmas}
\makeatletter
\newlabel{lemma-Artin-Rees}{{00IN}{source}{}{sourcepointer-1}{}}
\newlabel{lemma-Artin-Rees@cref}{{[originallemma][0][]00IN}{[1][0][]source}{}{}{}}
\makeatother
\crefname{originallemma}{original lemma}{original lemmas}
\Crefname{originallemma}{Original lemma}{Original lemmas}
\makeatletter
\newlabel{lemma-differ-finite-chi}{{00KB}{source}{}{sourcepointer-2}{}}
\newlabel{lemma-differ-finite-chi@cref}{{[originallemma][0][]00KB}{[1][0][]source}{}{}{}}
\makeatother
\crefname{originallemma}{original lemma}{original lemmas}
\Crefname{originallemma}{Original lemma}{Original lemmas}
\makeatletter
\newlabel{lemma-hilbert-ses}{{00K6}{source}{}{sourcepointer-3}{}}
\newlabel{lemma-hilbert-ses@cref}{{[originallemma][0][]00K6}{[1][0][]source}{}{}{}}
\makeatother
\title{The Hilbert-Samuel function in a short exact sequence, and degrees of the associated polynomials}
\author{The Stacks Project Authors}
\date{stacks-project:algebra.tex:local-snapshot-20261012; GNU FDL 1.2 or later}
\begin{document}
\maketitle
\noindent\textit{Editorial note.} All delivered mathematical text is the author's own, copied from the pinned Stacks Project chapter file algebra.tex; no statement or proof step was rewritten, reordered or removed.\par\medskip
\noindent\textit{Editorial note.} Preamble provenance, as required for this source: the frozen snapshot stages the chapter's own preamble file beside it (preamble.tex, 217 lines, 8116 bytes, the same shared Stacks-wide preamble the published chow and etale-cohomology cards use); the first line of the pinned chapter file is \textbackslash{}input\{preamble\}, so no preamble text lives inside the chapter itself. The card takes that file's lines 8-18 (the verbatim package and the three \textbackslash{}comment-based environments reference/slogan/history at lines 9-12, section 13.2, together with the Xy-pic setup) and its lines 154 to the runtime end of the file (the native theorem declarations and mathematical macros). Not one character was changed, no package was added, no macro was defined, and the excerpt serves only as preamble, never as a span coordinate.\par\medskip
\noindent\textit{Editorial note.} The same-chapter pointers are located at their \textbackslash{}label lines inside the pinned chapter file; the cross-chapter pointers are located in the project's own tag list staged in the snapshot (sources/stacks-chow/tags.txt), at the verbatim lines carrying their canonical tags.\par\medskip
\noindent\textit{Editorial note.} The counted proof ends by leaving the degenerate cases in which one of the modules vanishes to the reader; that sentence is kept verbatim inside the span and the main line is written out.\par\medskip
\noindent\textit{Editorial note.} The counted proof refers to same-chapter statements by label; they are recorded as navigation stubs pointed at their \textbackslash{}label lines.\par\medskip
\noindent\textit{Allowed standard prerequisites (editorial).}\begin{itemize}
\item Comparison of Hilbert-Samuel functions for a finite-length difference. Invoked in the counted proof at lines 14459-14462. Controls the difference of the functions for two modules differing by a finite length quotient.
\item Additivity of the Hilbert-Samuel function in a short exact sequence up to a bounded error. Invoked in the counted proof (lines 14462-14470). Gives the degree estimate required in part (1).
\end{itemize}\medskip
\ifdefined\hypertarget\hypertarget{sourcepointer-1}{}\fi
\noindent\textit{Source reference.} Original lemma 00IN: sources/stacks-algebra/pinned-original-source.tex, source lines 12404--12404. Navigation pointer for \textbackslash{}ref\{lemma-Artin-Rees\} used inside the counted span (line 14478); the argument that consumes it is retained in the span.\par\medskip
\ifdefined\hypertarget\hypertarget{sourcepointer-2}{}\fi
\noindent\textit{Source reference.} Original lemma 00KB: sources/stacks-algebra/pinned-original-source.tex, source lines 14430--14430. Navigation pointer for \textbackslash{}ref\{lemma-differ-finite-chi\} used inside the counted span (line 14460); the argument that consumes it is retained in the span.\par\medskip
\ifdefined\hypertarget\hypertarget{sourcepointer-3}{}\fi
\noindent\textit{Source reference.} Original lemma 00K6: sources/stacks-algebra/pinned-original-source.tex, source lines 14305--14305. Navigation pointer for \textbackslash{}ref\{lemma-hilbert-ses\} used inside the counted span (line 14459); the argument that consumes it is retained in the span.\par\medskip
\begin{lemma}
\label{lemma-hilbert-ses-chi}
Let $R$ be a Noetherian local ring. Let $I \subset R$ be an ideal of
definition. Let $0 \to M' \to M \to M'' \to 0$ be a short exact sequence
of finite $R$-modules. Then
\begin{enumerate}
\item if $M'$ does not have finite length, then
$\chi_{I, M} - \chi_{I, M''} - \chi_{I, M'}$
is a numerical polynomial of degree $<$ the degree of
$\chi_{I, M'}$,
\item $\max\{ \deg(\chi_{I, M'}), \deg(\chi_{I, M''}) \} = \deg(\chi_{I, M})$,
and
\item $\max\{d(M'), d(M'')\} = d(M)$,
\end{enumerate}
\end{lemma}

\begin{proof}
We first prove (1). Let $N \subset M'$ be as in Lemma \ref{lemma-hilbert-ses}.
By Lemma \ref{lemma-differ-finite-chi} the numerical polynomial
$\chi_{I, M'} - \chi_{I, N}$ has degree $<$ the common degree of
$\chi_{I, M'}$ and $\chi_{I, N}$. By Lemma \ref{lemma-hilbert-ses}
the difference
$$
\chi_{I, M}(n) - \chi_{I, M''}(n) - \chi_{I, N}(n - c)
$$
is constant for $n \gg 0$. By elementary calculus the difference
$\chi_{I, N}(n) - \chi_{I, N}(n - c)$ has degree $<$ the degree of
$\chi_{I, N}$ which is bigger than zero (see above). Putting everything
together we obtain (1).

\medskip\noindent
Note that the leading coefficients of $\chi_{I, M'}$ and $\chi_{I, M''}$ are
nonnegative. Thus the degree of $\chi_{I, M'} + \chi_{I, M''}$ is equal
to the maximum of the degrees. Thus if $M'$ does not have finite
length, then (2) follows from (1). If $M'$ does have finite length, then
$I^nM \to I^nM''$ is an isomorphism for all $n \gg 0$ by Artin-Rees
(Lemma \ref{lemma-Artin-Rees}). Thus $M/I^nM \to M''/I^nM''$ is a
surjection with kernel $M'$ for $n \gg 0$ and we see that
$\chi_{I, M}(n) - \chi_{I, M''}(n) = \text{length}(M')$
for all $n \gg 0$. Thus (2) holds in this case also.

\medskip\noindent
Proof of (3). This follows from (2) except if one of $M$, $M'$, or $M''$
is zero. We omit the proof in these special cases.
\end{proof}

\end{document}
